\documentclass[11pt]{article}

\usepackage[margin=1in]{geometry}
\usepackage{amsmath, amssymb}
\usepackage{enumitem}
\usepackage{booktabs}

\setlength{\parindent}{0pt}
\setlength{\parskip}{0.8em}

\begin{document}

\begin{center}
{\Large \textbf{Midterm Practice}}\\
\vspace{0.3em}
\end{center}

\section*{Easy}

\begin{enumerate}[label=\textbf{\arabic*.}]

\item Fund A accumulates at an annual nominal interest rate of 12\%, convertible monthly.

Fund B accumulates with an annual force of interest, $\delta_t=\dfrac{t}{6}$ for all $t$.

At time $t=0$, 1 is deposited in each fund.

$T$, in years, is the time that the two funds are equal, $T>0$.

Determine $T$.

\begin{itemize}
\item A. $12\ln(1.01)$
\item B. $12\bigl[\ln(1.12)-\ln(1.01)\bigr]$
\item C. $12(1.12)$
\item D. $144\ln(1.01)$
\item E. $144\ln(1.12)$
\end{itemize}

\medskip
\hrule
\medskip

\item The prices for four 1,000 face amount zero-coupon bonds are as follows:

\[
\begin{array}{|c|c|}
\hline
\textbf{Price} & \textbf{Term (in years)}\\
\hline
943.40 & 1\\
\hline
747.26 & 5\\
\hline
558.39 & 10\\
\hline
311.80 & 20\\
\hline
\end{array}
\]

Determine which of the following statements describes the yield curve underlying these prices.

\begin{itemize}
\item A. The yield curve is increasing with constant slope.
\item B. The yield curve is flat.
\item C. The yield curve is inverted with constant slope.
\item D. The yield curve is concave upward.
\item E. The yield curve is concave downward.
\end{itemize}

\medskip
\hrule
\medskip

\item The accumulation function is given by
\[
a(t)=(1+t)^{1/2}.
\]

What is the force of interest, $\delta_t$, at time, $t=1$?

\begin{itemize}
\item A. $1/2$
\item B. $1/3$
\item C. $1/4$
\item D. $2/3$
\item E. $3/4$
\end{itemize}

\end{enumerate}

\section*{Medium}

\begin{enumerate}[label=\textbf{\arabic*.}]

\item A perpetuity pays at a continuous rate of $500t$ per year, where $t$ represents the time in years since the start of the perpetuity. The annual force of interest is $\delta_t = 0.02$.

Calculate the present value of the perpetuity.

\begin{itemize}
\item A. 500,000
\item B. 1,000,000
\item C. 1,250,000
\item D. 1,500,000
\item E. 2,500,000
\end{itemize}

\medskip
\hrule
\medskip

\item The present value of a company’s liabilities is 5,000. The Macaulay duration and Macaulay convexity of the company’s liabilities are 12 and 195.

The company creates an investment portfolio consisting of investments in two of the following zero-coupon bonds:

\begin{enumerate}
\item[(i)] 5-year bond
\item[(ii)] 10-year bond
\item[(iii)] 20-year bond
\end{enumerate}

The position satisfies the requirements for Redington immunization. The annual effective yield on each of the bonds is 8%.

Determine which of the following portfolios is created by the company.

\begin{itemize}
\item A. Invest 1,000 in the 5-year bond and 4,000 in the 10-year bond
\item B. Invest 2,000 in the 5-year bond and 3,000 in the 20-year bond
\item C. Invest 2,666.67 in the 5-year bond and 2,333.33 in the 20-year bond
\item D. Invest 3,000 in the 10-year bond and 2,000 in the 20-year bond
\item E. Invest 4,000 in the 10-year bond and 1,000 in the 20-year bond
\end{itemize}

\medskip
\hrule
\medskip

\item A bank makes a loan of 20,000 at an interest rate of $i$.

The loan will be repaid with level payments at the end of each year for 20 years.

When the bank receives each payment, it reinvests at a rate of 5%.

At the end of the 20 year period, the bank calculates the annual effective return over the loan period to be 6.5%.

What is $i$, the original interest rate on the loan?

\begin{itemize}
\item A. 4.42%
\item B. 6.57%
\item C. 8.62%
\item D. 9.04%
\item E. 10.6%
\end{itemize}

\medskip
\hrule
\medskip

\item A common stock pays a constant dividend at the end of each year into perpetuity.

Using an annual effective interest rate of 8%, calculate the Macaulay convexity of the stock.

\begin{itemize}
\item A. 324
\item B. 325
\item C. 351
\item D. 352
\item E. 378
\end{itemize}

\medskip
\hrule
\medskip

\item On January 1, 1987, three 100 par value bonds with 6% annual coupons will mature at the end of 1, 2, and 3 years, respectively. The redemption value of each bond is 100.

You are given that the prices for these bonds on January 1, 1987 are:

\[
\begin{array}{|c|c|}
\hline
\textbf{Maturity Date} & \textbf{Price} \\
\hline
\text{December } 31, 1987 & 101.92 \\
\hline
\text{December } 31, 1988 & 102.84 \\
\hline
\text{December } 31, 1989 & 105.51 \\
\hline
\end{array}
\]

These prices are based on an interest rate of $i$ in 1987, $j$ in 1988, and $k$ in 1989.

Determine $j$.

\begin{itemize}
\item A. 2%
\item B. 3%
\item C. 4%
\item D. 5%
\item E. 6%
\end{itemize}

\medskip
\hrule
\medskip

\item A 36-year annuity-immediate with annual payments of \$4 has the same present value as an 18-year annuity-immediate with annual payments of \$5.

In how many years does money double at rate of interest, $i$?

\begin{itemize}
\item A. 9
\item B. 12
\item C. 18
\item D. 27
\item E. 36
\end{itemize}

\end{enumerate}

\newpage
\section*{Hard}

\begin{enumerate}[label=\textbf{\arabic*.}]

\item Bill purchases an annuity at a price of 10,000. The annuity makes payments of 500 at the beginning of every 6 months for 20 years.

The payments are reinvested in a fund which earns interest at an annual effective rate $i$.

Interest payments are received every 6 months and reinvested at a nominal rate of 6%, convertible semiannually.

Bill realizes an overall effective annual yield of 7% on his original investment over the 20-year period.

Calculate $i$.

\begin{itemize}
\item A. 5.90%
\item B. 6.05%
\item C. 6.20%
\item D. 6.35%
\item E. 6.50%
\end{itemize}

\medskip
\hrule
\medskip

\item At an effective annual interest rate $i$, you are given:

i. the present value of an annuity-immediate with annual payments of 1 for $n$ years is 40; and

ii. the present value of an annuity-immediate with annual payments of 1 for $3n$ years is 70.

Calculate the accumulated value of an annuity-immediate with annual payments of 1 for $2n$ years.

\begin{itemize}
\item A. 240
\item B. 243
\item C. 260
\item D. 268
\item E. 280
\end{itemize}

\medskip
\hrule
\medskip

\item Consider two annuities that have the same yield and make annual payments.

The first annuity is a level $n$-year annuity-due, and its Macaulay duration is 3.

The second annuity is a level $n$-year annuity-immediate, and its Macaulay duration is $X$.

Calculate $X$.

\begin{itemize}
\item A. 3.00
\item B. 3.25
\item C. 3.50
\item D. 4.00
\item E. 4.50
\end{itemize}

\medskip
\hrule
\medskip

\item A perpetuity with annual payments is payable beginning 10 years from now. The first payment is 50. Each annual payment thereafter is increased by 10 until a payment of 150 is reached. Subsequent payments remain level at 150.

This perpetuity is purchased by means of 10 annual premiums, with the first premium of $P$ due immediately. Each premium after the first is 105 percent of the preceding one.

The annual effective interest rates are 5\% during the first 9 years and 3\% thereafter.

Calculate $P$.

\begin{itemize}
\item A. 281
\item B. 286
\item C. 291
\item D. 296
\item E. 301
\end{itemize}

\end{enumerate}


\newpage

Solutions

\textbf{Easy}

\begin{enumerate}
    \item D
    \item B
    \item C
\end{enumerate}

\textbf{Medium}

\begin{enumerate}
    \item C
    \item C
    \item C
    \item C
    \item D
    \item A
\end{enumerate}

\textbf{Hard}

\begin{enumerate}
    \item B
    \item A
    \item D
    \item C
\end{enumerate}



\end{document}